% R. Brown's translations of two letters to L. Breen, with his note and comments (cote 134-2): archive cote 134-2, site page 128 (PDF page 129).
% translation of the letter of 17-19 July 1975. Draft compared: ps1.tex, lines 1380-1404.
% AI-transcribed by Claude Opus 5.5 (claude-opus-5-5, adaptive thinking, effort high), Message Batches API, batch brown_full (msgbatch_01Md3i5yK1PxkpJn5GHfEExE), 2026-11-13; not yet checked by eye.
\scanpage{134-2/128}
% NOTE: a white patch on the scan, after "from" in the first title line, hides part of the line
\uline{Translation of an extract from}

% NOTE: a white patch on the scan, before "a Letter" in the second title line, hides part of the line
\uline{a Letter from A. Grothendieck to L. Breen}

\uline{17-19 July, 1975}

I am happy to finish by receiving an echo to my long letter
and even a beginning to a constructive approach to a theorem of
the type I envisaged.  The construction which you propose for the
notion of a non-strict n-category, and of the functor nerve, has
certainly the merit of existing, and of being a first precise
approach, but otherwise can be  subject to some evident criticism:
it is very technical, hardly intuitive (yet at the level of
1-cat, etc, for 2-cat everything  is so clear ``you just follow
your nose ...'')  And finally the absence of a definition of a
functor sending (semi-)simplicial sets to n-groupoids.  This
functor corresponds to a geometric intuition so clear that a theory
which does not include it seems to me to be useless.  Perhaps
in trying to write down (like a sort of list of Christmas presents!)
in a complete and explicit enough way the notions which one would
like to have at ones disposal, and the relations (functor,
equivalence, etc) one will arrive finally at a kind of axiomatic
description sufficiently complete to be able to give the key to
an explicit \uline{ad} \uline{hoc} construction, or should permit maybe to enunciate
and prove a theorem of existence and uniqueness for a theory of
the required type.

Otherwise, not having understood the idea of Segal in your
last letter (which I have generously sent to Illusie ...), I do
not see how you define the Picard n-categories - but this matters
little.  As far as ``strict'' Picard n-categories are concerned,
what I ask of them is that they finally form an (n+1)-category

\mspage{\hand{129}}
